\lecture{3}{На хор}

\section{На хор}

\subsection{Ключевые слова и примеры}
\subsubsection{\texttt{static}}
\subsubsection{\texttt{explicit}}
\subsubsection{\texttt{mutable}}
\subsubsection{\texttt{friend}}

\subsection{Продвинутые шаблоны}
\subsubsection{Три вида шаблонных параметров}
\subsubsection{Перегрузка шаблонных функций}
\subsubsection{Специализация шаблонных классов и функций}
\subsubsection{Шаблонные аргументы по умолчанию}
\subsubsection{Шаблонные параметры шаблонов}
\subsubsection{Проблема зависимых имён (оба вида)}

\subsection{Variadic templates и type traits}
\subsubsection{Синтаксис и примеры}
\subsubsection{Оператор \texttt{sizeof...}}
\subsubsection{Базовые type-traits: \texttt{is\_same}, \texttt{remove\_const}}
\subsubsection{Type-traits с variadic: \texttt{is\_same\_many}}

\subsection{Продвинутое наследование}
\subsubsection{Private, public, protected наследование}
\subsubsection{Видимость/доступность полей и методов}
\subsubsection{Влияние \texttt{friend} при наследовании}
\subsubsection{Приведение типов при всех типах наследования}
\subsubsection{Множественное наследование}
\subsubsection{Проблема ромба}
\subsubsection{Расположение объектов в памяти}
\subsubsection{Ошибки: \textquotedblleft ambiguous base\textquotedblright, \textquotedblleft warning: inaccessible base class\textquotedblright}

\subsection{Виртуальные приколы}
\subsubsection{Аргументы по умолчанию в виртуальных функциях}
\subsubsection{Вызов чисто виртуальных функций}
\subsubsection{Приведение типов в рантайме: \texttt{dynamic\_cast}}
\subsubsection{\texttt{override} и \texttt{final}}

\subsection{Продвинутые лямбды}
\subsubsection{Список захвата}
\subsubsection{Захват с инициализацией}
\subsubsection{Захват полей класса и \texttt{this}}
\subsubsection{Захват по умолчанию}
\subsubsection{Лямбды как объекты}
\subsubsection{\texttt{std::function}: идея и пример использования}

\subsection{Продвинутые исключения}
\subsubsection{Проблема исключений в конструкторах/деструкторах}
\subsubsection{Идиома RAII}
\subsubsection{Спецификация исключений и \texttt{noexcept}}
\subsubsection{Оператор \texttt{noexcept}}
\subsubsection{Стандартная иерархия исключений}
\subsubsection{Кастомные исключения}
\subsubsection{Приведение типов внутри \texttt{catch}}

\subsection{Итераторы}
\subsubsection{\texttt{std::iterator\_traits}}
\subsubsection{\texttt{std::advance}: реализация}
\subsubsection{const и reverse итераторы}
\subsubsection{\texttt{std::copy}, \texttt{std::copy\_if}}
\subsubsection{output итераторы и адаптеры}
\subsubsection{stream итераторы}
\subsubsection{Итератор для вектора: псевдокод}

\subsection{Инвалидация ссылок/итераторов и устройство контейнеров}
\subsubsection{Инвалидация в стандартных контейнерах}
\subsubsection{Подробное устройство \texttt{std::deque}, \texttt{std::list}, \texttt{std::unordered\_map}}

\subsection{Продвинутая память}
\subsubsection{Перегрузка операторов \texttt{new} и \texttt{delete}}
\subsubsection{Placement new}
\subsubsection{Идея аллокаторов}
\subsubsection{Аллокаторы в \texttt{vector}}
\subsubsection{\texttt{std::allocator} и его методы}
\subsubsection{\texttt{std::allocator\_traits}}
\subsubsection{Реализация \texttt{resize} вектора}

\subsection{Категории значений}
\subsubsection{Наивное определение lvalue/rvalue: контрпримеры}
\subsubsection{Правильное определение lvalue/rvalue}
\subsubsection{rvalue-ссылки и правила их инициализации}
\subsubsection{Правила присваивания rvalue-ссылок}
\subsubsection{RVO и copy-elision: примеры}
\subsubsection{Reference qualifiers}

\subsection{Универсальные ссылки и \texttt{std::move}}
\subsubsection{Правила вывода шаблонов для универсальных ссылок}
\subsubsection{Реализация \texttt{std::move}}

\subsection{Perfect forwarding}
\subsubsection{Описание проблемы}
\subsubsection{\texttt{std::forward}: пример использования}
\subsubsection{Правильная реализация \texttt{emplace\_back} и \texttt{resize} вектора}

\subsection{\texttt{shared\_ptr} и сопутствующее}
\subsubsection{Идея реализации \texttt{shared\_ptr} (с \texttt{weak\_ptr})}
\subsubsection{\texttt{make\_shared}, \texttt{make\_unique}: пользование и реализация}

\subsection{SFINAE и проверка наличия метода}
\subsubsection{Идея SFINAE, простейший пример}
\subsubsection{\texttt{std::enable\_if}: реализация и использование}
\subsubsection{Проверка наличия метода в классе: идея реализации}

\subsection{requires и концепты}
\subsubsection{Ограничения через SFINAE-характеристики}
\subsubsection{Сложные ограничения}
\subsubsection{Концепты: примеры и сокращённый синтаксис}
\subsubsection{Вложение концептов: простой пример}
